% !TEX encoding = UTF-8
% !TEX Root = Main.tex
\section{Initialization Technique}
\label{sec:dnmt}

The initialization technique ADMIN \cite{liu2020deeptransformer} we will apply here reformulates Eq.~\ref{eq:postnorm} as:
\vspace{-2mm}
\begin{equation}
    \mathbf{x_{i}}=f_{LN}(\mathbf{x_{i-1}}\cdot \mathbf{\omega_i} + f_{i}(\mathbf{x_{i-1}}))
    \label{eq:admin}
\end{equation}
where $\mathbf{\omega_i}$ is a constant vector that is element-wise multiplied to $\mathbf{x_{i-1}}$ in order to balance the contribution against $f_{i}(\mathbf{x_{i-1}})$. 
The observation is that in addition to vanishing gradients, the unequal magnitudes in the two terms 
$\mathbf{x_{i-1}}$ and $f_{i}(\mathbf{x_{i-1})}$ is the main cause of instability in training.
Refer to \cite{liu2020deeptransformer} for theoretical details.\footnote{Note that paper presents results of 18-layer Transformers on the WMT'14 En-De, which we also use here. Our contribution is a more comprehensive evaluation.} 

ADMIN initialization involves two phases:
At the {\bf Profiling} phase, we randomly initialize the model parameters using default initialization, set $\mathbf{\omega_i}=1$, and perform one step forward pass in order to compute the output variance of the residual branch $Var[f(\mathbf{x_{i-1})}]$ at each layer.\footnote{We estimate the variance with one batch of 8k tokens.} In the {\bf Training} phase, we fix $\mathbf{\omega_i}~=~\sqrt{\sum_{j < i}{Var[f(\mathbf{x_{j-1}})]}}$, and then train the model using standard back-propagation. 
After training finishes, $\mathbf{\omega_i}$ can be folded back into the model parameters to recover the standard Transformer architecture. 
This simple initialization method is effective in ensuring that training does not diverge, even in deep networks. 